\textbf{Phần 1. Trắc nghiệm lựa chọn}
\setcounterpageref{ex}{0}

\begin{ex} %[3X1V74]
Trong một hộp có 5 viên bi màu xanh và không có viên bi nào màu khác. Nhắm mắt lấy ra 1 viên bi. Sự kiện nào sau đây là 'Chắc chắn' xảy ra?
\choice
{Lấy được 1 viên bi màu đỏ}
{\True Lấy được 1 viên bi màu xanh}
{Lấy được 1 viên bi màu vàng}
{Không lấy được viên bi nào}
\loigiai{
Vì trong hộp chỉ toàn bi xanh nên chắc chắn sẽ lấy được viên bi màu xanh.
\textbf{Đáp án của câu hỏi là B}
}
\end{ex}

\begin{ex} %[3X1V74]
Gieo một con xúc xắc chuẩn có 6 mặt ghi từ 1 đến 6 chấm. Sự kiện nào sau đây là 'Không thể' xảy ra?
\choice
{Mặt xuất hiện có 6 chấm}
{Mặt xuất hiện có 1 chấm}
{Mặt xuất hiện có số chấm là số chẵn}
{\True Mặt xuất hiện có 7 chấm}
\loigiai{
Xúc xắc chỉ có từ 1 đến 6 chấm nên không thể xuất hiện mặt 7 chấm.
\textbf{Đáp án của câu hỏi là D}
}
\end{ex}

\begin{ex} %[3X1V74]
Trong túi có 3 chiếc kẹo dâu và 4 chiếc kẹo cam. Nhắm mắt lấy ra 1 chiếc kẹo. Khả năng xảy ra là
\choice
{\True Có thể lấy được kẹo dâu hoặc kẹo cam}
{Chắc chắn lấy được kẹo dâu}
{Chắc chắn lấy được kẹo cam}
{Không thể lấy được kẹo dâu}
\loigiai{
Trong túi có cả kẹo dâu và kẹo cam nên có thể lấy được kẹo dâu hoặc kẹo cam.
\textbf{Đáp án của câu hỏi là A}
}
\end{ex}

\begin{ex} %[3X1V74]
Gieo một đồng xu có hai mặt: mặt Sấp và mặt Ngửa. Khả năng xuất hiện mặt Sấp là
\choice
{Chắc chắn xảy ra}
{Không thể xảy ra}
{\True Có thể xảy ra}
{Luôn luôn xảy ra}
\loigiai{
Đồng xu có hai khả năng Sấp hoặc Ngửa nên khả năng xuất hiện mặt Sấp là có thể xảy ra.
\textbf{Đáp án của câu hỏi là C}
}
\end{ex}

\begin{ex} %[3X1V74]
Trong một hộp có 4 chiếc bút mực màu đỏ và 4 chiếc bút mực màu xanh. Lấy ra cùng một lúc 5 chiếc bút. Sự kiện nào sau đây là 'Chắc chắn' xảy ra?
\choice
{Tất cả 5 chiếc đều màu đỏ}
{\True Có ít nhất 1 chiếc bút màu đỏ và 1 chiếc bút màu xanh}
{Tất cả 5 chiếc đều màu xanh}
{Không có chiếc bút màu xanh nào}
\loigiai{
Vì chỉ có 4 bút đỏ nên lấy 5 chiếc chắc chắn có ít nhất 1 bút xanh; và chỉ có 4 bút xanh nên chắc chắn có ít nhất 1 bút đỏ.
\textbf{Đáp án của câu hỏi là B}
}
\end{ex}

\begin{ex} %[3X1V74]
Ngày mai Mặt Trời mọc ở hướng nào là sự kiện 'Chắc chắn' xảy ra?
\choice
{Mặt Trời mọc ở hướng Tây}
{Mặt Trời mọc ở hướng Bắc}
{Mặt Trời mọc ở hướng Nam}
{\True Mặt Trời mọc ở hướng Đông}
\loigiai{
Theo quy luật tự nhiên, Mặt Trời chắc chắn mọc ở hướng Đông.
\textbf{Đáp án của câu hỏi là D}
}
\end{ex}

\begin{ex} %[3X1V74]
Trong một bình có 10 viên bi đỏ và 1 viên bi vàng. Bạn Nam lấy ngẫu nhiên 1 viên bi. Nhận xét nào sau đây là đúng nhất?
\choice
{\True Khả năng lấy được bi đỏ cao hơn khả năng lấy được bi vàng}
{Chắc chắn lấy được bi đỏ}
{Không thể lấy được bi vàng}
{Khả năng lấy được bi đỏ và bi vàng là như nhau}
\loigiai{
Vì số bi đỏ nhiều hơn số bi vàng nên khả năng lấy được bi đỏ cao hơn.
\textbf{Đáp án của câu hỏi là A}
}
\end{ex}

\begin{ex} %[3X1V74]
Hôm nay là thứ Sáu. Sự kiện 'Ngày mai là thứ Bảy' là sự kiện
\choice
{Có thể xảy ra}
{Không thể xảy ra}
{\True Chắc chắn xảy ra}
{Tùy thuộc thời tiết}
\loigiai{
Theo quy luật tuần hoàn các ngày trong tuần, sau thứ Sáu chắc chắn là thứ Bảy.
\textbf{Đáp án của câu hỏi là C}
}
\end{ex}

\begin{ex} %[3X1V74]
Một con cá vàng sống trong chậu nước. Sự kiện 'Con cá có thể trèo lên cây chuối' là
\choice
{Chắc chắn xảy ra}
{\True Không thể xảy ra}
{Có thể xảy ra}
{Chưa chắc chắn}
\loigiai{
Cá không thể trèo cây, đây là sự kiện không thể xảy ra.
\textbf{Đáp án của câu hỏi là B}
}
\end{ex}

\begin{ex} %[3X1V74]
Một hộp có 3 quả bóng xanh và 2 quả bóng đỏ. Nam lấy ra một số quả bóng. Để CHẮC CHẮN có ít nhất 1 quả bóng màu đỏ thì Nam phải lấy ra ít nhất bao nhiêu quả bóng?
\choice
{$2$ quả bóng}
{$3$ quả bóng}
{$5$ quả bóng}
{\True $4$ quả bóng}
\loigiai{
Trường hợp xấu nhất lấy phải 3 quả bóng xanh, cần lấy thêm 1 quả nữa: $3 + 1 = 4$ quả bóng.
\textbf{Đáp án của câu hỏi là D}
}
\end{ex}

\begin{ex} %[3X1V74]
Gieo đồng thời hai con xúc xắc. Tổng số chấm xuất hiện lớn nhất có thể là
\choice
{\True $12$}
{$6$}
{$10$}
{$14$}
\loigiai{
Số chấm lớn nhất ở mỗi xúc xắc là 6: $6 + 6 = 12$ chấm.
\textbf{Đáp án của câu hỏi là A}
}
\end{ex}

\begin{ex} %[3X1V74]
Bốc ngẫu nhiên một tấm thẻ trong 10 tấm thẻ ghi số từ 1 đến 10. Sự kiện 'Rút được tấm thẻ ghi số 11' là
\choice
{Chắc chắn xảy ra}
{Có thể xảy ra}
{\True Không thể xảy ra}
{Luôn xảy ra}
\loigiai{
Các thẻ chỉ ghi từ 1 đến 10 nên không thể rút được số 11.
\textbf{Đáp án của câu hỏi là C}
}
\end{ex}

\textbf{Phần 2. Trắc nghiệm đúng sai}
\setcounterpageref{ex}{0}

\begin{ex} %[3X1V74]
Xét tính đúng sai của các sự kiện ngẫu nhiên trong cuộc sống:
\choiceTF
{\True Một sự kiện chắc chắn xảy ra là sự kiện luôn luôn xảy ra trong mọi trường hợp.}
{Một sự kiện có thể xảy ra là sự kiện không bao giờ xảy ra.}
{\True Một sự kiện không thể xảy ra là sự kiện không bao giờ xảy ra trong bất kì trường hợp nào.}
{Gieo 1 con xúc xắc chắc chắn xuất hiện mặt 6 chấm.}
\loigiai{
\begin{itemize}
    \item Ý a) Đúng: Định nghĩa.
    \item Ý b) Sai: Sai vì sự kiện có thể xảy ra tức là có khả năng xảy ra.
    \item Ý c) Đúng: Định nghĩa.
    \item Ý d) Sai: Sai vì có thể xuất hiện các mặt từ 1 đến 5.
\end{itemize}
\begin{itemize}
    \item a) \textbf{Đúng}.
    \item b) \textbf{Sai}.
    \item c) \textbf{Đúng}.
    \item d) \textbf{Sai}.
\end{itemize}
}
\end{ex}

\begin{ex} %[3X1V74]
Trong một hộp kín có 8 viên bi trắng và 2 viên bi đen:
\choiceTF
{\True Lấy ngẫu nhiên 1 viên bi thì có thể lấy được bi trắng hoặc bi đen.}
{Chắc chắn lấy được viên bi màu đen khi chỉ lấy 1 viên.}
{\True Nếu lấy ra 9 viên bi thì chắc chắn có ít nhất 1 viên bi màu đen.}
{Không thể lấy được viên bi màu trắng.}
\loigiai{
\begin{itemize}
    \item Ý a) Đúng: Có cả hai loại bi.
    \item Ý b) Sai: Sai vì có thể lấy phải bi trắng.
    \item Ý c) Đúng: Vì bi trắng chỉ có 8 viên.
    \item Ý d) Sai: Sai vì trong hộp có 8 bi trắng.
\end{itemize}
\begin{itemize}
    \item a) \textbf{Đúng}.
    \item b) \textbf{Sai}.
    \item c) \textbf{Đúng}.
    \item d) \textbf{Sai}.
\end{itemize}
}
\end{ex}

\textbf{Phần 3. Trả lời ngắn}
\setcounterpageref{ex}{0}

\begin{ex} %[3X1V74]
Trong hộp có 5 viên bi đỏ và 3 viên bi vàng. Cần lấy ra ít nhất bao nhiêu viên bi để chắc chắn có ít nhất 1 viên bi màu vàng?
\shortans{6}
\loigiai{
\textbf{Đáp số:} $6$\\
Trường hợp xấu nhất lấy 5 bi đỏ, lấy thêm 1 viên nữa: $5 + 1 = 6$ viên.
}
\end{ex}

\begin{ex} %[3X1V74]
Gieo một con xúc xắc. Tổng số mặt có số chấm là số chẵn là bao nhiêu mặt?
\shortans{3}
\loigiai{
\textbf{Đáp số:} $3$\\
Các mặt chẵn: 2, 4, 6 (có 3 mặt).
}
\end{ex}

\begin{ex} %[3X1V74]
Trong túi có 10 viên kẹo sô-cô-la và không có loại kẹo nào khác. Lấy ra 1 chiếc kẹo. Có bao nhiêu khả năng về màu sắc/loại kẹo có thể nhận được?
\shortans{1}
\loigiai{
\textbf{Đáp số:} $1$\\
Chỉ có duy nhất 1 loại kẹo sô-cô-la.
}
\end{ex}

\begin{ex} %[3X1V74]
Gieo 1 con xúc xắc 6 mặt. Số chấm nhỏ nhất có thể xuất hiện là bao nhiêu?
\shortans{1}
\loigiai{
\textbf{Đáp số:} $1$\\
Mặt 1 chấm là nhỏ nhất.
}
\end{ex}

\begin{ex} %[3X1V74]
Có 4 quả táo và 5 quả lê trong giỏ. Hỏi cần lấy ra ít nhất bao nhiêu quả để chắc chắn có ít nhất 1 quả táo?
\shortans{6}
\loigiai{
\textbf{Đáp số:} $6$\\
Trường hợp xấu nhất lấy cả 5 quả lê, thêm 1 quả nữa: $5 + 1 = 6$ quả.
}
\end{ex}

\begin{ex} %[3X1V74]
Một tuần lễ có bao nhiêu ngày chắc chắn có tên bắt đầu bằng chữ 'Thứ'?
\shortans{6}
\loigiai{
\textbf{Đáp số:} $6$\\
Thứ Hai, Thứ Ba, Thứ Tư, Thứ Năm, Thứ Sáu, Thứ Bảy (6 ngày, trừ Chủ nhật).
}
\end{ex}
